\Needspace{12\baselineskip}
\originalchapter{丝状聚合物与分子马达}
\originalsection{生长、缩短和动态失稳}
细胞骨架并非静止支架。聚合末端可以添加或失去单体，整个链长是随机过程。固定增长率 $a$、缩短率 $b$ 的最简模型若允许长度从 $1$ 开始，在长度 $1$ 处禁止进一步缩短，则稳态可归一化的条件是 $a<b$，几何分布 $p_\ell=(1-a/b)(a/b)^{\ell-1}$。当 $a>b$ 时概率向大长度迁移，不存在此半无限状态空间上的归一化稳态。

微管还在快速生长态与快速缩短态间切换。令灾变率 $\alpha$、恢复率 $\gamma$，增长速度 $v_+>0$、缩短速度的大小 $v_->0$。长期处于两态的比例为 $\pi_+=\gamma/(\alpha+\gamma)$、$\pi_-=\alpha/(\alpha+\gamma)$，有效漂移
\[
v=\frac{\gamma v_+-\alpha v_-}{\alpha+\gamma}.
\]
速度涨落的关联按 $\e^{-(\alpha+\gamma)t}$ 衰减，方差幅度是 $\pi_+\pi_-(v_++v_-)^2$。将速度关联对时间积分得到附加扩散 $D_{\rm switch}=\alpha\gamma(v_++v_-)^2/(\alpha+\gamma)^3$。这说明即使每种状态内部没有随机位移，随机切换也会产生长期扩散。

\originalsection{马达的内部循环}
分子马达消耗ATP等化学燃料沿丝前进。一轮步进距离为 $d$，若要通过 $m$ 个独立等速率 $u$ 的内部阶段，等待时间是 $m$ 个指数变量之和，均值 $m/u$、方差 $m/u^2$。更新过程给出 $v=du/m$、$D=d^2u/(2m^2)$，所以随机度 $r=2D/(dv)=1/m$。低随机度支持多个限速阶段，但有分支、逆步和非指数等待时，不能直接把 $1/r$ 当作精确内部状态数。

有逆向跃迁的循环须满足局部详细平衡。一轮前后速率乘积比通常含化学驱动力与负载功之差：\[\ln\frac{\prod_i u_i}{\prod_i w_i}=\beta(\Delta\mu-Fd).\]零平均速度的停滞力可反映这一比值，然而停滞也可能仍有无效耗能循环。平均速度为零不等于整个系统热平衡。

\originalsection{趋化与随机游动}
大肠杆菌的奔跑和翻滚使方向不断更新。每段奔跑时长 $t_r$、速度 $v$，随后的翻滚时长 $t_t$，新方向均匀抽取。在 $d_s$ 个空间维度中，不同奔跑位移交叉平均为零，因此长期
\[
D=\frac{v^2\avg{t_r^2}}{2d_s(\avg{t_r}+\avg{t_t})}.
\]
这里出现二阶矩，而非仅平均奔跑时间，因为长奔跑对位移平方贡献更大。若沿营养梯度的方向平均奔跑时间稍长，每轮的位移期望不再为零，便出现漂移。将 $
\tau_r(\hat n)=\tau_0+g\hat n\cdot\hat z$ 代入 $\avg{\hat n(\hat n\cdot\hat z)}=\hat z/d_s$，得到 $v_{\rm drift}=vg/[d_s(\tau_0+\tau_t)]$。局部营养感知的时间记忆是更细的机制，不能由此简化模型直接证明。

\originalsection{微管长度分布与边界再注入}
增长/缩短态长度密度满足
\begin{align*}
\partial_tp_+&=-\alpha p_++\gamma p_- -v_+\partial_xp_+,\\
\partial_tp_-&=\alpha p_+-\gamma p_- +v_-\partial_xp_-.
\end{align*}
此处 $v_-$ 为正的缩短速率大小。总长度流为 $J=v_+p_+-v_-p_-$。若到零处立即重新进入增长态，边界要求零净流；若微管消失或有独立成核，则另有吸收/源项，不能沿用相同稳态。

零流稳态令 $p_+/p_-=v_-/v_+$，代回得到 $p_\pm=C_\pm\e^{-\lambda x}$，
\[
\lambda=\frac{\alpha v_- -\gamma v_+}{v_+v_-},\qquad
\avg x=\lambda^{-1}.
\]
要求 $\lambda>0$，等价于无限直线上的长期漂移为负。值得注意：有限长度边界稳态的增长/缩短态比例是速度比，而不是无限空间中仅内部切换的 $\gamma/\alpha$；到零再注入改变了时间占据。把两种比值混用会破坏零流条件。

\originalsection{随机切换产生扩散}
在远离边界的长时间极限，内部指示变量 $I_+(t)$ 的均值为 $\pi_+$，协方差 $\pi_+\pi_-\e^{-(\alpha+\gamma)|t|}$。速度 $V=(v_++v_-)I_+-v_-$。位移涨落方差
\[
\operatorname{Var}X(t)=2\int_0^t(t-s)\avg{\delta V(s)\delta V(0)}\dd s
\simeq2t\frac{\alpha\gamma(v_++v_-)^2}{(\alpha+\gamma)^3}.
\]
这个推导保留内部状态弛豫修正；只把每态概率代成瞬时平衡再相加通常会漏掉二阶扩散项。下面的马达模型也有同样问题。

\originalsection{帽的片段化与灾变时间}
若GTP帽只在末端加长、在后边界水解，帽长像漂移扩散过程，正漂移会无限长大，不能解释更复杂的灾变观测。源讲义引入每单位长度水解率 $r$：长度 $x$ 的帽被任意内部水解事件截短，损失率 $rx$，来自更长帽的生成率为 $r\int_x^\infty p(y)\dd y$。
\[
\partial_tp=-v\partial_xp+D_c\partial_x^2p-rxp+r\int_x^\infty p(y)\dd y.
\]
定义尾积分 $P(x)=\int_x^\infty p(y)\dd y$，对上式积分得到 $\partial_tP=-vP'+D_cP''-rxP$。这一步把非局部项变成局部方程，而不是忽略片段化。

在富含单体、$v>0$ 且先忽略 $D_c$ 时，稳态 $P_*=\e^{-rx^2/(2v)}$，所以 $p_*=(rx/v)\e^{-rx^2/(2v)}$，平均帽长 $\sqrt{\pi v/(2r)}$。小扩散使概率穿过零帽长，首阶灾变率约 $D_cr/v$；它是对漂移主导稳态的扰动结果，并非所有参数下精确。

换到贫单体溶液，$v<0$。从足够长初帽开始，能够在时刻 $t$ 前导致灾变的水解事件构成 $(x,\tau)$ 平面面积 $|v|t^2/2$ 的三角域。Poisson无事件概率为 $\e^{-r|v|t^2/2}$，故
\[
p_{\rm cat}(t)=r|v|t\e^{-r|v|t^2/2},\qquad
\avg t_{\rm cat}=\sqrt{\frac\pi{2r|v|}}.
\]
它不依赖原帽长，前提是初帽足以覆盖这段事件域；短帽仍会受初始边界影响。这是利用随机几何解释等待时间分布的例子。

\originalsection{两态马达的准确长时统计}
记轨道步长 $d$，内部状态T、D，T到同格D率为 $d_0$、到右格D率为 $r$；D到同格T率为 $u$、到左格T率为 $l$。内部总弛豫率 $S=u+d_0+r+l$。倾斜矩阵的非对角元为 $u+l\e^{-kd}$ 与 $d_0+r\e^{kd}$，最大特征值 $\Lambda(k)$ 是位移累积量生成率，其方程
\[
\Lambda^2+S\Lambda-ru(\e^{kd}-1)-ld_0(\e^{-kd}-1)=0.
\]
在 $k=0$ 求导，得到
\[
v=\frac{d(ru-ld_0)}S,\qquad
D=\frac{d^2(ru+ld_0)}{2S}-\frac{d^2(ru-ld_0)^2}{S^3}.
\]
最后负项描述等待时间与内部状态的相关。原讲义先作瞬时局部平衡后给出的扩散表达一般不等于此精确长时 $D$，所选PS8题解已采用倾斜生成元复核，并在QA记录此差异；漂移式则一致。

\originalsection{化学功、停滞与力的不同定义}
一周期净化学自由能 $\Delta\mu$ 与反负载功 $Fd$ 决定循环亲和力 $\mathcal A=\beta(\Delta\mu-Fd)=\ln(ru/(ld_0))$。紧耦合单循环的停滞条件为 $\mathcal A=0$，给 $F_s=\Delta\mu/d$。若有无位移耗ATP的旁路，停滞时仍可消耗燃料，不能从零速度推导零耗散。

定义 $F_E=\kb Tv/D$ 是将马达数据套进平衡粒子的Einstein关系得到的比较尺度。主动马达一般不在平衡线性响应区，此值不自动等于真实外力或最大负载；应明确采用准确的 $D$ 或近似的 $D$。功率为 $P=Fv$，停滞点功率为零，最大功率通常位于更小负载。热力学效率 $Fv/(J_{\rm ATP}\Delta\mu)$ 还需独立知道耗燃料通量，不能只由 $v,D$ 唯一推出。


\originalsection{细胞游动的低Re极限与scallop定理}
分子马达把化学功变为内部运动，但整个细胞怎样在水中前进，还取决于流体响应。细胞尺度a、速度U、流体黏度 $\eta$ 与密度 $\rho$ 给 $Re=\rho Ua/\eta$。当它小、形变周期长于黏性扩散时间 $\tau_\nu=\rho a^2/\eta$，忽略流体惯性得到瞬时Stokes方程
\[
 -\nabla p+\eta\nabla^2\mathbf u=0,\qquad \nabla\cdot\mathbf u=0.
\]
附加无滑移运动边界与远场条件后，速度场和牵引力线性依赖边界速度。游动者不靠惯性滑行：形变停下、驱动力撤去后，它不能按宏观船只的惯性继续远行。布朗运动和方向扩散也会限制保持方向的时间。

用q描述体形，U描述沿固定轴的整体平移，线性阻力可写 $F=-R_{UU}(q)U-R_{Uq}(q)\dot q$。无外力游动要求 $F=0$，因此 $U=A(q)\dot q$，$A=-R_{Uq}/R_{UU}$。若一个往返周期只沿同一个形状参数区间先开再合，
\[
 \Delta X=\int A(q)\dd q=0.
\]
返程的积分恰好抵消去程；快开慢合不会改变这一结论，因为在瞬时线性、无记忆模型中速度变化与持续时间相抵。这是scallop定理的最小形式。包括旋转的完整刚体运动也具有相应的时间反演性：沿同一形状路径逆向返回会撤销前进过程。原课程明确以低Re细胞运动和此定理为主题，此处为编者方程化说明。

两种独立形变 $q_1,q_2$ 可以在形状空间围成环。仍只考虑沿固定轴的平移，小振幅闭路有 $\Delta X\simeq(\partial_1A_2-\partial_2A_1)\mathcal S$，$\mathcal S$为有向面积；交换动作次序会反转位移。这使“先伸后弯”和“先弯后伸”成为不同周期。对于带整体转动的三维运动，还应使用包含旋转的连接与非交换项，不能机械套用这一标量面积式。非Newton流体记忆、显著惯性或外场变化则可破坏上面的假设，须另建模型。

旋转鞭毛也不是一个铰链的往返：旋转角持续绕圈，螺旋手性使转动与平移耦合。以一个所选转动模式的有效阻力示意，
\[
 \begin{pmatrix}F\\N\end{pmatrix}
 =-\begin{pmatrix}A&B\\B&C\end{pmatrix}
 \begin{pmatrix}U\\\Omega\end{pmatrix},\qquad
 U=-\frac BA\Omega\quad(F=0).
\]
正耗散要求 $A>0,AC-B^2>0$，故维持该模式所需功率为 $(C-B^2/A)\Omega^2>0$。非螺旋对称结构可有 $B=0$，此时转动不能在该轴向产生推进。真实自由细胞的马达施加内力矩，细胞体反向旋转；完整问题还须把细胞体与鞭毛分别列入力和力矩平衡。这一简化阻力矩阵解释手性耦合，不冒充完整菌体流场计算。

\originalsection{跑动、翻滚与时间上的趋化比较}
细胞太小时，对一瞬间横跨自身的浓度差进行测量很困难；跑动过程中可比较先前与当前的化学信号，调整转向概率。跑动不是始终朝正确方向；降低有利方向的终止率就可产生统计偏向，同时方向扩散和翻滚重新探索其他方向。细胞爬行还涉及actin重塑与基底接触力，不能用无外力液体游动的同一个边界条件解释。

一个明确的教学模型只保留一维速度 $+v_0,-v_0$，向负、向正翻转的率分别为 $\alpha_+,\alpha_-$。稳态方向概率为 $p_+=\alpha_-/(\alpha_++\alpha_-)$，故
\[
 v_{\rm drift}=v_0\frac{\alpha_--\alpha_+}{\alpha_++\alpha_-},\qquad
 D_{\rm eff}=\frac{4v_0^2\alpha_+\alpha_-}{(\alpha_++\alpha_-)^3}.
\]
后式来自速度涨落方差 $4v_0^2\alpha_+\alpha_-/(\alpha_++\alpha_-)^2$ 乘相关时间 $1/(\alpha_++\alpha_-)$。对称转向时漂移为零而扩散仍为 $v_0^2/(2\alpha)$；有利方向的 $\alpha_+$下降后出现正漂移。这里忽略三维转角、信号适应和旋转Brown运动，只示范一个内部状态规则如何变成宏观输运，和前面两态马达的长时统计相接。
